% ═══════════════════════════════════════════════════════════════
% SL-02d-ML: Stundenleistung MUSTERLÖSUNG
% Spanisch Klasse 7 - Sequenz 2: Mi familia
% ═══════════════════════════════════════════════════════════════

\documentclass[11pt, a4paper]{article}

\newif\ifswmodus
\swmodusfalse

\usepackage[utf8]{inputenc}
\usepackage[T1]{fontenc}
\usepackage[ngerman]{babel}
\usepackage{helvet}
\renewcommand{\familydefault}{\sfdefault}

\usepackage[margin=2cm]{geometry}
\usepackage{enumitem}
\usepackage{tabularx}
\usepackage{booktabs}

\usepackage{xcolor}
\usepackage{tcolorbox}
\tcbuselibrary{skins}

\usepackage{fancyhdr}
\usepackage{lastpage}
\usepackage{needspace}

\definecolor{spanisch-color}{HTML}{c41e3a}
\definecolor{grau-color}{HTML}{888888}
\definecolor{hellgrau-color}{HTML}{f5f5f5}
\definecolor{loesung-color}{HTML}{c62828}

\definecolor{sw-dark}{HTML}{000000}
\definecolor{sw-gray}{HTML}{666666}
\definecolor{sw-white}{HTML}{ffffff}

\ifswmodus
    \colorlet{spanisch}{sw-dark}
    \colorlet{grau}{sw-gray}
    \colorlet{hellgrau}{sw-white}
    \colorlet{loesung}{sw-dark}
\else
    \colorlet{spanisch}{spanisch-color}
    \colorlet{grau}{grau-color}
    \colorlet{hellgrau}{hellgrau-color}
    \colorlet{loesung}{loesung-color}
\fi

\newcommand{\fachfarbe}{spanisch}

\newcommand{\thema}{Stundenleistung: Mi familia}
\newcommand{\fach}{Spanisch}
\newcommand{\klasse}{7}
\newcommand{\maxpunkte}{20}
\newcommand{\zeit}{25 Minuten}
\newcommand{\datum}{\today}

\pagestyle{fancy}
\fancyhf{}
\renewcommand{\headrulewidth}{0.4pt}
\renewcommand{\footrulewidth}{0.4pt}

\lhead{\fach{} -- Klasse \klasse}
\chead{\textcolor{loesung}{\textbf{SL-02d MUSTERLÖSUNG}}}
\rhead{\datum}

\lfoot{\zeit}
\cfoot{}
\rfoot{Seite \thepage\ von \pageref{LastPage}}

\newcommand{\punktebox}[1]{\hfill\fbox{\small \_/#1}}
\newcommand{\luecke}[1]{\rule{#1}{0.4pt}}
\newcommand{\lsg}[1]{\textcolor{loesung}{\textbf{#1}}}

\newtcolorbox{infobox}{
    colback=hellgrau,
    colframe=\fachfarbe,
    boxrule=1pt,
    arc=0pt,
    left=8pt, right=8pt, top=8pt, bottom=8pt
}

\newenvironment{aufgabe}[1]{%
    \needspace{5\baselineskip}%
    \nopagebreak[4]%
    \vspace{10pt}
    \noindent\textbf{\textcolor{\fachfarbe}{Aufgabe #1}}
    \nopagebreak[4]%
    \vspace{5pt}
    \nopagebreak[4]%
    \begin{list}{}{\leftmargin=0pt}
    \item[]
}{%
    \end{list}
}

\newcommand{\es}[1]{\textcolor{\fachfarbe}{\textbf{#1}}}

\begin{document}

\begin{center}
    {\Large\bfseries\textcolor{\fachfarbe}{\thema}}\\[0.3em]
    {\large\textcolor{loesung}{\textbf{--- MUSTERLÖSUNG ---}}}
\end{center}

\vspace{5pt}

\noindent
\begin{tabularx}{\textwidth}{|X|X|X|}
\hline
\textbf{Name:} \lsg{Musterschüler/in} &
\textbf{Klasse:} \klasse &
\textbf{Datum:} \datum \\
\hline
\end{tabularx}

\vspace{10pt}

\begin{infobox}
\textbf{Zeit:} \zeit{} | \textbf{Punkte:} \maxpunkte{} | \textbf{Hilfsmittel:} keine
\end{infobox}

\vspace{15pt}

\begin{aufgabe}{1 -- Familienmitglieder zuordnen}
Verbinde die spanischen Familienbegriffe mit der deutschen Bedeutung. \punktebox{4}
\end{aufgabe}

\begin{center}
\begin{tabular}{rcl}
\es{la abuela} & \lsg{---} & der Bruder \lsg{(= la abuela $\rightarrow$ die Großmutter)} \\[0.8em]
\es{el hermano} & \lsg{---} & die Cousine \lsg{(= el hermano $\rightarrow$ der Bruder)} \\[0.8em]
\es{la prima} & \lsg{---} & die Großmutter \lsg{(= la prima $\rightarrow$ die Cousine)} \\[0.8em]
\es{el tío} & \lsg{---} & der Onkel \lsg{(= el tío $\rightarrow$ der Onkel)} \\
\end{tabular}
\end{center}

\textit{\textcolor{grau}{Bewertung: 1 Punkt pro richtige Zuordnung}}

\vspace{15pt}

\begin{aufgabe}{2 -- Das Verb \es{tener} (haben) einsetzen}
Ergänze die richtige Form von \es{tener}. \punktebox{4}
\end{aufgabe}

\noindent
a) Yo \lsg{tengo} dos hermanas. \hfill (ich habe)

\vspace{0.8em}

\noindent
b) María \lsg{tiene} un perro. \hfill (María hat)

\vspace{0.8em}

\noindent
c) Nosotros \lsg{tenemos} una casa grande. \hfill (wir haben)

\vspace{0.8em}

\noindent
d) ¿Tú \lsg{tienes} primos? \hfill (hast du?)

\textit{\textcolor{grau}{Bewertung: 1 Punkt pro richtige Form}}

\vspace{15pt}

\begin{aufgabe}{3 -- Possessivbegleiter und Zahlen}
Ergänze die Lücken mit dem passenden Possessivbegleiter oder der Zahl. \punktebox{6}
\end{aufgabe}

\noindent
a) \lsg{Mi} madre tiene 45 años. \hfill (meine Mutter)

\vspace{0.8em}

\noindent
b) ¿Cuántos años tiene \lsg{tu} hermano? \hfill (dein Bruder)

\vspace{0.8em}

\noindent
c) \lsg{Nuestros} abuelos viven en España. \hfill (unsere Großeltern)

\vspace{0.8em}

\noindent
d) Mi hermana tiene \lsg{doce} años. \hfill (zwölf)

\vspace{0.8em}

\noindent
e) \lsg{Sus} primos tienen \lsg{quince} años. \hfill (seine Cousins / fünfzehn)

\textit{\textcolor{grau}{Bewertung: 1 Punkt pro richtige Lücke}}

\vspace{15pt}

\begin{aufgabe}{4 -- Dialog vervollständigen}
Ergänze den Dialog über die Familie sinnvoll auf Spanisch. \punktebox{6}
\end{aufgabe}

\noindent
\textbf{A:} ¡Hola! ¿Tienes hermanos?

\vspace{0.8em}

\noindent
\textbf{B:} \lsg{Sí, tengo un hermano. / Sí, tengo dos hermanas. / No, no tengo hermanos.} \hfill (2P)

\vspace{0.8em}

\noindent
\textbf{A:} ¿Cómo se llama tu madre?

\vspace{0.8em}

\noindent
\textbf{B:} \lsg{Mi madre se llama [Name]. / Se llama María.} \hfill (2P)

\vspace{0.8em}

\noindent
\textbf{A:} ¿Cuántos años tiene tu padre?

\vspace{0.8em}

\noindent
\textbf{B:} \lsg{Mi padre tiene [Zahl] años. / Tiene cuarenta y cinco años.} \hfill (2P)

\textit{\textcolor{grau}{Bewertung: 2 Punkte pro sinnvolle, grammatisch korrekte Antwort}}

\vspace{25pt}
\hrule
\vspace{10pt}

\noindent
\begin{tabularx}{\textwidth}{|X|c|c|c|c||c|}
\hline
\textbf{Aufgabe} & 1 & 2 & 3 & 4 & \textbf{Gesamt} \\
\hline
\textbf{max. Punkte} & 4 & 4 & 6 & 6 & \textbf{20} \\
\hline
\textbf{erreicht} & \lsg{4} & \lsg{4} & \lsg{6} & \lsg{6} & \lsg{20} \\
\hline
\end{tabularx}

\vspace{10pt}

\begin{flushright}
\textbf{Note:} \lsg{1}
\end{flushright}

\vspace{1em}

\textbf{Notenschlüssel:}

\begin{tabular}{|c|c|c|c|c|c|c|}
\hline
\textbf{Note} & 1 & 2 & 3 & 4 & 5 & 6 \\
\hline
\textbf{ab Punkte} & 19 & 16 & 12 & 8 & 4 & <4 \\
\hline
\end{tabular}

\vspace{1em}

\textbf{Hinweise zur Bewertung:}
\begin{itemize}
    \item Aufgabe 1: Zuordnung muss eindeutig sein (Striche, Zahlen, Pfeile)
    \item Aufgabe 2: Auch ``tengo'' ohne Akzent akzeptieren
    \item Aufgabe 3: ``Nuestros'' statt ``nuestro'' wegen Plural (abuelos)
    \item Aufgabe 4: Teilpunkte möglich (1P bei kleinen Fehlern)
    \item Generell: Akzentfehler bei sonst korrekter Schreibung tolerieren
\end{itemize}

\end{document}
